\documentclass[aspectratio=169,12pt]{beamer}

% ============================================================
%  PACKAGES
% ============================================================
\usepackage[T1]{fontenc}
\usepackage[utf8]{inputenc}
\usepackage[spanish]{babel}
\usepackage{amsmath,amssymb,bm}
\usepackage{tikz}
\usetikzlibrary{arrows.meta,positioning,calc,backgrounds}
\usepackage{tcolorbox}
\tcbuselibrary{skins,breakable}
\usepackage{listings}
\usepackage{xcolor}
\usepackage{booktabs}
\usepackage{colortbl}

% ============================================================
%  COLORS
% ============================================================
\definecolor{darkbg}{HTML}{1A1A2E}
\definecolor{accent}{HTML}{E94560}
\definecolor{accentteal}{HTML}{0F9B8E}
\definecolor{accentamber}{HTML}{F5A623}
\definecolor{accentviolet}{HTML}{7B68EE}
\definecolor{textlight}{HTML}{EAEAEA}
\definecolor{textgray}{HTML}{9E9E9E}
\definecolor{lightbg}{HTML}{F7F7F7}
\definecolor{coralbox}{HTML}{FF6B6B}

% ============================================================
%  BEAMER SETUP
% ============================================================
\setbeamercolor{background canvas}{bg=lightbg}
\setbeamercolor{normal text}{fg=darkbg}
\setbeamertemplate{navigation symbols}{}
\setbeamertemplate{footline}{%
  \hfill\textcolor{textgray}{\tiny\texttt{mbastidaso@unal.edu.co}\quad\insertframenumber/\inserttotalframenumber\quad}%
  \vspace{4pt}
}

% ============================================================
%  DARK SLIDE ENVIRONMENT
% ============================================================
\newenvironment{darkslide}{%
  \setbeamercolor{background canvas}{bg=darkbg}%
  \setbeamercolor{normal text}{fg=textlight}%
}{%
  \setbeamercolor{background canvas}{bg=lightbg}%
  \setbeamercolor{normal text}{fg=darkbg}%
}

% ============================================================
%  TCOLORBOX ENVIRONMENTS
% ============================================================
\tcbset{mybase/.style={
  enhanced, breakable, arc=3pt,
  boxrule=0.6pt, left=6pt, right=6pt, top=5pt, bottom=5pt,
  fonttitle=\small\bfseries, fontupper=\small
}}

\newtcolorbox{defbox}[1]{mybase,
  colback=accentviolet!8, colframe=accentviolet!60,
  title={#1}}

\newtcolorbox{resultbox}[1]{mybase,
  colback=accentteal!8, colframe=accentteal!60,
  title={#1}}

\newtcolorbox{obsbox}[1]{mybase,
  colback=accentamber!10, colframe=accentamber!70,
  title={#1}}

\newtcolorbox{taskbox}[1]{mybase,
  colback=coralbox!8, colframe=coralbox!60,
  title={#1}}

% ============================================================
%  LISTINGS
% ============================================================
\lstset{
  language=Python,
  basicstyle=\ttfamily\tiny,
  keywordstyle=\color{accentviolet}\bfseries,
  commentstyle=\color{textgray},
  stringstyle=\color{accentteal},
  backgroundcolor=\color{darkbg!5},
  frame=single, framerule=0.4pt,
  breaklines=true, showstringspaces=false,
  morekeywords={jax,jit,grad,vmap,jnp}
}

% ============================================================
%  TIKZ STYLES
% ============================================================
\tikzset{
  neuron/.style={circle, draw=#1, fill=#1!15, minimum size=9mm, thick, font=\small},
  conn/.style={-{Stealth[length=2mm]}, gray!60, line width=0.5pt},
  thickconn/.style={-{Stealth[length=2.5mm]}, #1, line width=1.2pt},
  graphnode/.style={circle, draw=accentviolet, fill=accentviolet!12,
                    minimum size=8mm, thick, font=\footnotesize},
  graphop/.style={rectangle, rounded corners=2pt, draw=accentteal,
                  fill=accentteal!12, minimum width=12mm, minimum height=7mm,
                  thick, font=\footnotesize},
}

% ============================================================
%  DOCUMENT
% ============================================================
\begin{document}

% ============================================================
%  TÍTULO
% ============================================================
\begin{darkslide}
\begin{frame}[plain]
\begin{tikzpicture}[remember picture, overlay]
  \fill[accent,opacity=0.15] (current page.north east) ++(-1.5,-1.5) circle (4cm);
  \fill[accentteal,opacity=0.10] (current page.south west) ++(2,2) circle (3cm);
  \fill[accentviolet,opacity=0.08] (current page.north west) ++(1,-2) circle (2.5cm);
\end{tikzpicture}
\vspace{1.5cm}
\begin{center}
  {\small\textcolor{textgray}{Programación Científica · 2026-1}}\\[16pt]
  {\LARGE\textcolor{textlight}{\textbf{Redes Neuronales}}}\\[6pt]
  {\large\textcolor{accent}{Otra vez: Aproximar funciones}}\\[20pt]
  {\normalsize\textcolor{textgray}{Universidad Nacional de Colombia}}\\[4pt]
  {\small\textcolor{textgray}{\texttt{mbastidaso@unal.edu.co}}}
\end{center}
\end{frame}
\end{darkslide}

% ============================================================
%  PREGUNTA DE APERTURA
% ============================================================
\begin{darkslide}
\begin{frame}[plain]
\begin{tikzpicture}[remember picture, overlay]
  \fill[accentteal,opacity=0.12] (current page.south east) ++(-2,2) circle (3.5cm);
\end{tikzpicture}
\vspace{2.5cm}
\begin{center}
  {\large\textcolor{textgray}{La pregunta de hoy:}}\\[18pt]
  {\LARGE\textcolor{textlight}{\textbf{¿Cómo aproximamos}}}\\[6pt]
  {\LARGE\textcolor{accent}{\textbf{funciones muy complicadas?}}}\\[30pt]
  {\small\textcolor{textgray}{(dejen esta pregunta abierta — la responderemos al final)}}
\end{center}
\end{frame}
\end{darkslide}

% ============================================================
%  SECCIÓN 1: HILO DEL CURSO
% ============================================================
\begin{darkslide}
\begin{frame}[plain]
\begin{tikzpicture}[remember picture, overlay]
  \fill[accentviolet,opacity=0.12] (current page.north east) ++(-2,-2) circle (3cm);
  \fill[accent,opacity=0.10] (current page.south west) ++(1,1) circle (2.5cm);
\end{tikzpicture}
\vspace{2cm}
\begin{center}
  {\Large\textcolor{textlight}{\textbf{Parte 1}}}\\[10pt]
  {\large\textcolor{accentteal}{El hilo del curso}}
\end{center}
\end{frame}
\end{darkslide}

\begin{frame}{El hilo del curso}

\vspace{6pt}
\begin{resultbox}{El mismo problema, con más libertad}
En todos los métodos anteriores buscamos:
\[
  \hat{f}(x) = \sum_{i=1}^n c_i\,\phi_i(x)
\]
Las bases $\phi_i$ estaban \textbf{fijas}. En una red neuronal, los $\phi_i$ \textbf{también dependen de parámetros} que se optimizan.
\end{resultbox}
\end{frame}

\begin{frame}{¿Qué ganamos con bases aprendibles?}
\begin{columns}[T]
\column{0.48\textwidth}
\begin{obsbox}{Bases fijas}
\begin{itemize}\setlength\itemsep{3pt}
  \item Escogemos la forma de $\phi_i$ antes de ver los datos
  \item Funciona bien si tenemos intuición sobre la estructura
  \item Monomios: oscilación de Runge
  \item Splines: necesitamos nodos bien ubicados
\end{itemize}
\end{obsbox}

\column{0.48\textwidth}
\begin{resultbox}{Bases aprendibles}
\begin{itemize}\setlength\itemsep{3pt}
  \item Los $\phi_i$ se adaptan a los datos
  \item No necesitamos saber la estructura a priori
  \item El precio: necesitamos \textbf{optimizar} los parámetros
  \item ¿Cómo? → gradiente 
\end{itemize}
\end{resultbox}
\end{columns}

\vspace{8pt}
\begin{center}
\textcolor{accentviolet}{\textbf{Hoy:}} construimos la maquinaria mínima para entender cómo se optimiza.
\end{center}
\end{frame}

% ============================================================
%  SECCIÓN 2: RED NEURONAL MÍNIMA
% ============================================================
\begin{darkslide}
\begin{frame}[plain]
\begin{tikzpicture}[remember picture, overlay]
  \fill[accent,opacity=0.12] (current page.north west) ++(2,-2) circle (3cm);
  \fill[accentteal,opacity=0.10] (current page.south east) ++(-2,2) circle (3cm);
\end{tikzpicture}
\vspace{2cm}
\begin{center}
  {\Large\textcolor{textlight}{\textbf{Parte 2}}}\\[10pt]
  {\large\textcolor{accent}{La red neuronal mínima}}
\end{center}
\end{frame}
\end{darkslide}

\begin{frame}{Una neurona}
\begin{columns}[T]
\column{0.52\textwidth}
\begin{defbox}{Neurona (ML)}
Recibe entradas $x_1,\ldots,x_n$, calcula:
\[
  z = \sum_{j=1}^n w_j x_j + b, \qquad a = \sigma(z)
\]
\begin{itemize}
  \item $w_j$: pesos (\textit{weights})
  \item $b$: sesgo (\textit{bias})
  \item $\sigma$: función de activación (no lineal)
\end{itemize}
\end{defbox}

\column{0.44\textwidth}
\vspace{4pt}
\begin{center}
\begin{tikzpicture}[node distance=7mm and 12mm]
  \node[neuron=textgray] (x1) {$x_1$};
  \node[neuron=textgray, below=of x1] (x2) {$x_2$};
  \node[neuron=textgray, below=of x2] (x3) {$x_3$};
  \node[neuron=accentviolet, right=22mm of x2] (n) {$\sigma(z)$};
  \node[neuron=accentteal, right=12mm of n] (out) {$a$};

  \draw[conn] (x1) -- node[above,font=\small,sloped]{$w_1$} (n);
  \draw[conn] (x2) -- node[above,font=\small]{$w_2$} (n);
  \draw[conn] (x3) -- node[below,font=\small,sloped]{$w_3$} (n);
  \draw[thickconn=accentteal] (n) -- (out);

  \node[above=3mm of n, font=\small\color{textgray}] {$z = \bm{w}^\top\bm{x}+b$};
\end{tikzpicture}
\end{center}

\vspace{6pt}
\underline{{Activaciones comunes:}}\\

{\footnotesize
$\sigma(z) = \tanh(z)$\\
$\sigma(z) = \max(0,z)$ \textcolor{textgray}{(ReLU)}\\
$\sigma(z) = (1+e^{-z})^{-1}$ \textcolor{textgray}{(sigmoid)}
}
\end{columns}
\end{frame}

\begin{frame}{Red feedforward de una capa oculta}

\begin{defbox}{Arquitectura}
\[
  \hat{f}(x;\theta) = w^2\,\sigma(W^1 x + \bm{b}^1) + b_1^2
\]
\begin{itemize}
  \item Capa oculta: $\bm{h} = \sigma(W_1 x + \bm{b}_1) \in \mathbb{R}^m$
  \item Capa de salida: $\hat{f} = W_2\bm{h} + b_2 \in \mathbb{R}$
  \item Parámetros: $\theta = \{W_1, \bm{b}_1, W_2, b_2\}$
\end{itemize}
\end{defbox}

\begin{obsbox}{Dimensiones}
Si $x \in \mathbb{R}$ y tengo $m$ neuronas ocultas:\\
$W_1 \in \mathbb{R}^{m\times 1}$, $\bm{b}_1 \in \mathbb{R}^m$, $W_2 \in \mathbb{R}^{1\times m}$, $b_2 \in \mathbb{R}$\\
\textbf{Total:} $3m + 1$ parámetros.
\end{obsbox}

\end{frame}

\begin{frame}[fragile]{Dibujemos un grafo}
\begin{taskbox}{Ejercicio}
Considere la red con $m=2$ neuronas ocultas y entrada escalar $x$:
\[
  h_1 = \sigma(w_{1}^1x + b_1^1), \quad h_2 = \sigma(w_{2}^1x + b_2^1), \quad \hat{f} = w_1^2 h_1 + w_2^2 h_2 + b_1^2
\]
\textbf{Dibujen} el grafo computacional. \\[4pt]
¿Cuántos parámetros tiene esta red?
\end{taskbox}

\vspace{-0.7cm}
\begin{columns}[T]
\column{0.6\textwidth}
\begin{obsbox}{Variables intermedias}
$z_1 = w_{1}^1x + b_1$\\
$h_1 = \sigma(z_1)$\\
$z_2 = w_{2}^1x + b_2$\\
$h_2 = \sigma(z_2)$\\
$\hat{f} = w_1^2 h_1 + w_2^2 h_2 + b_1^2$
\end{obsbox}
\column{0.46\textwidth}

\end{columns}
\end{frame}

% ============================================================
%  SECCIÓN 3: LOSS FUNCTION
% ============================================================
\begin{darkslide}
\begin{frame}[plain]
\begin{tikzpicture}[remember picture, overlay]
  \fill[accentamber,opacity=0.12] (current page.north east) ++(-2,-2) circle (3cm);
  \fill[accentteal,opacity=0.08] (current page.south west) ++(1,1) circle (3.5cm);
\end{tikzpicture}
\vspace{2cm}
\begin{center}
  {\Large\textcolor{textlight}{\textbf{Parte 3}}}\\[10pt]
  {\large\textcolor{accentamber}{La función de pérdida}}
\end{center}
\end{frame}
\end{darkslide}

\begin{frame}{¿Cómo medimos el error?}
Queremos que $\hat{f}(x;\theta) \approx f(x)$ para todo $x \in [a,b]$.

\vspace{8pt}
\begin{defbox}{Loss function (pérdida L2)}
\[
  \mathcal{L}(\theta) = \int_a^b \bigl(f(x) - \hat{f}(x;\theta)\bigr)^2\,dx
\]
\end{defbox}

\vspace{6pt}
En la práctica, aproximamos la integral con $N$ puntos de muestra:

\begin{resultbox}{Pérdida discreta (MSE)}
\[
  \mathcal{L}(\theta) \approx \frac{1}{N}\sum_{i=1}^N \bigl(f(x_i) - \hat{f}(x_i;\theta)\bigr)^2
\]
Los $\{x_i\}$ se llaman \textbf{datos de entrenamiento}.
\end{resultbox}

\vspace{4pt}
\begin{obsbox}{Observación clave}
$\mathcal{L}$ es una función $\mathcal{L}: \mathbb{R}^p \to \mathbb{R}$ donde $p = |\theta|$.\\
Derivamos con respecto a $\theta$, \textbf{no} con respecto a $x$.
\end{obsbox}
\end{frame}

%\begin{frame}{Geometría de la pérdida}
%\begin{columns}[T]
%\column{0.55\textwidth}
%\vspace{4pt}
%Imaginen $\mathcal{L}$ como una superficie sobre el espacio de parámetros $\theta \in \mathbb{R}^p$:
%
%\vspace{8pt}
%\begin{center}
%\begin{tikzpicture}[scale=0.9]
%  % bowl shape using ellipses
%  \draw[accentviolet!30, line width=1pt] (0,0) ellipse (2.5 and 0.8);
%  \draw[accentviolet!50, line width=1pt] (0,0) ellipse (1.8 and 0.55);
%  \draw[accentviolet!70, line width=1pt] (0,0) ellipse (1.0 and 0.30);
%  \fill[accent] (0,0) circle (2pt) node[right,font=\tiny\color{accent}] {mínimo};
%  \draw[-{Stealth}, accentteal, thick] (1.8,0.55) -- (1.2,0.36)
%    node[midway, right, font=\tiny\color{accentteal}] {$-\nabla_\theta\mathcal{L}$};
%  \node[above=3mm, font=\small\color{textgray}] at (0,0.8) {curvas de nivel de $\mathcal{L}(\theta)$};
%\end{tikzpicture}
%\end{center}
%
%\column{0.42\textwidth}
%\vspace{4pt}
%\begin{defbox}{Descenso de gradiente}
%Repetir hasta converger:
%\[
%  \theta \leftarrow \theta - \eta\,\nabla_\theta\mathcal{L}(\theta)
%\]
%$\eta > 0$: tasa de aprendizaje.
%\end{defbox}
%
%\vspace{6pt}
%\begin{obsbox}{El problema central}
%Necesitamos $\nabla_\theta\mathcal{L}$.\\[4pt]
%$\mathcal{L}$ depende de $\theta$ \textit{a través} de la red $\hat{f}(x;\theta)$.\\[4pt]
%¿Cómo calculamos ese gradiente?
%\end{obsbox}
%\end{columns}
%\end{frame}

% ============================================================
%  SECCIÓN: TEORÍA DEL GRADIENTE DESCENDENTE
% ============================================================
\begin{darkslide}
	\begin{frame}[plain]
		\begin{tikzpicture}[remember picture, overlay]
			\fill[accentamber,opacity=0.12] (current page.north west) ++(2,-2) circle (3.5cm);
			\fill[accentviolet,opacity=0.10] (current page.south east) ++(-2,2) circle (3cm);
		\end{tikzpicture}
		\vspace{2cm}
		\begin{center}
			{\Large\textcolor{textlight}{\textbf{Teoría}}}\\[10pt]
			{\large\textcolor{accentamber}{¿cómo optimizamos?}}
		\end{center}
	\end{frame}
\end{darkslide}

% --- Caso lineal: MSE como forma cuadrática ---
\begin{frame}{La pérdida MSE es un funcional cuadrático}
	Sea $\hat{f}(x;\theta) = \theta^\top \phi(x)$ una red \textbf{lineal} en $\theta$
	(e.g.\ regresión lineal, red sin activación).
	Con $N$ datos $\{(x_i, y_i)\}$:
	\[
	\mathcal{L}(\theta)
	= \frac{1}{N}\sum_{i=1}^N (y_i - \theta^\top\vec{\phi}(x_i))^2
	= \frac{1}{N}\|y - A\theta\|^2
	\]
	donde $y = (y_1,\ldots,y_N)^\top$ y $A_{ij} = \phi_j(x_i)$.
	
	\vspace{6pt}
%	\begin{resultbox}{Forma cuadrática explícita}
%		\[
%		\mathcal{L}(\theta)
%		= \frac{1}{N}\bigl(\theta^\top A^\top A\,\theta - 2\,y^\top A\,\theta + \|y\|^2\bigr)
%		= \frac{1}{N}\bigl(\theta^\top Q\,\theta - 2\,\ell^\top\theta + c\bigr)
%		\]
%		con $Q = A^\top A \in \mathbb{R}^{p\times p}$ semidefinida positiva,
%		$\ell = A^\top y$, $c = \|y\|^2$.
%	\end{resultbox}
	
	\begin{obsbox}{Conexión con mínimos cuadrados}
		$\nabla_\theta \mathcal{L} = 0$ da las ecuaciones normales $Q\theta^* = \ell$,
		exactamente el sistema de mínimos cuadrados que ya conocen.
		El mínimo global existe y es único si $A$ tiene rango completo.
	\end{obsbox}
\end{frame}

% --- Gradiente del funcional cuadrático ---
\begin{frame}{Gradiente de la pérdida cuadrática}
	\begin{resultbox}{Gradiente explícito (caso lineal)}
		\[
		\nabla_\theta \mathcal{L}(\theta)
		= \frac{2}{N} A^\top(A\theta - y)
		= \frac{2}{N} A^\top(\hat{y} - y)
		\]
		donde $\hat{y} = A\theta$ son las predicciones actuales.
	\end{resultbox}
	
	\vspace{6pt}
	\begin{obsbox}{¿Por qué el gradiente apunta al mínimo?}
		Para cualquier $\theta$, el vector $-\nabla_\theta\mathcal{L}(\theta)$ apunta
		en la dirección de \textbf{mayor descenso local} de $\mathcal{L}$.
		Para una forma cuadrática convexa, este descenso es global.
	\end{obsbox}
	
\end{frame}


% --- Generalización: caso no lineal ---
\begin{darkslide}
	\begin{frame}[plain]
		\begin{tikzpicture}[remember picture, overlay]
			\fill[accent,opacity=0.12] (current page.north east) ++(-2,-2) circle (3cm);
			\fill[accentteal,opacity=0.10] (current page.south west) ++(1,1) circle (3cm);
		\end{tikzpicture}
		\vspace{2cm}
		\begin{center}
			{\Large\textcolor{textlight}{\textbf{Generalización}}}\\[10pt]
			{\large\textcolor{accent}{Redes no lineales: $\hat{f}(x;\theta)$ arbitraria}}
		\end{center}
	\end{frame}
\end{darkslide}

% --- Aproximación cuadrática local ---
\begin{frame}{Aproximación cuadrática local (caso no lineal)}
	Para una red no lineal, $\mathcal{L}(\theta)$ puede ser no convexa.
	Pero localmente, alrededor de $\theta_k$, aproximamos por Taylor de orden 2:
	\[
	\mathcal{L}(\theta_k + \delta)
	\approx \mathcal{L}(\theta_k)
	+ \nabla_\theta\mathcal{L}(\theta_k)^\top \delta
	+ \tfrac{1}{2}\,\delta^\top H_k\,\delta,
	\]
	donde $H_k = \nabla^2_\theta\mathcal{L}(\theta_k)$ es la \textbf{matriz Hessiana}.
	
	\vspace{6pt}
	\begin{obsbox}{Mínimo local de la aproximación cuadrática}
		Si $H_k \succ 0$, el mínimo de la aproximación es:
		\[
		\delta^* = -H_k^{-1}\,\nabla_\theta\mathcal{L}(\theta_k).
		\]
		Gradiente descendente usa $\delta = -\eta\,\nabla_\theta\mathcal{L}$,
		es decir, aproxima $H_k \approx \frac{1}{\eta}I$.
	\end{obsbox}
	
%	\begin{resultbox}{Interpretación}
%		Cada paso de gradiente descendente minimiza exactamente
%		la aproximación cuadrática con Hessiana \textbf{esférica} $\frac{1}{\eta}I$.
%		Funciona bien cuando $\eta < 1/L$ donde $L = \|\nabla^2\mathcal{L}\|_2$
%		es la constante de Lipschitz del gradiente.
%	\end{resultbox}
\end{frame}

% --- Convergencia no lineal ---
\begin{frame}{Convergencia en el caso general}
	\begin{defbox}{Función $L$-suave}
		$\mathcal{L}$ es tal que $\|\nabla\mathcal{L}(\theta) - \nabla\mathcal{L}(\theta')\| \leq L\|\theta-\theta'\|$
		para todo $\theta, \theta'$, entonces $\nabla^2\mathcal{L}(\theta) \preceq LI$.
	\end{defbox}
	
	\begin{resultbox}{Teorema (descenso suficiente, caso $L$-suave)}
		Si $\mathcal{L}$ es $L$-suave y $\eta \leq 1/L$, entonces:
		\[
		\mathcal{L}(\theta_{k+1}) \leq \mathcal{L}(\theta_k) - \frac{\eta}{2}\|\nabla_\theta\mathcal{L}(\theta_k)\|^2.
		\]
	\end{resultbox}
	
%	\begin{proof}
%		Por $L$-suavidad: $\mathcal{L}(\theta_{k+1}) \leq \mathcal{L}(\theta_k)
%		+ \nabla\mathcal{L}_k^\top(\theta_{k+1}-\theta_k) + \frac{L}{2}\|\theta_{k+1}-\theta_k\|^2$.
%		Sustituyendo $\theta_{k+1}-\theta_k = -\eta\nabla\mathcal{L}_k$ y $\eta \leq 1/L$:
%		$\mathcal{L}(\theta_{k+1}) \leq \mathcal{L}(\theta_k) - \eta\|\nabla\mathcal{L}_k\|^2
%		+ \frac{\eta^2 L}{2}\|\nabla\mathcal{L}_k\|^2
%		\leq \mathcal{L}(\theta_k) - \frac{\eta}{2}\|\nabla\mathcal{L}_k\|^2$.
%	\end{proof}
	
	\begin{obsbox}{Consecuencia}
		$\mathcal{L}(\theta_k) \to$ mínimo local si $\nabla\mathcal{L}(\theta_k) \neq 0$
		en cada iteración. La convergencia a un \textbf{mínimo global} requiere convexidad.
	\end{obsbox}
\end{frame}

% --- Resumen de garantías ---
%\begin{frame}{Resumen: garantías de convergencia}
%	\begin{center}
%		\begin{tabular}{lccl}
%			\toprule
%			Escenario & Paso $\eta$ & Tasa & Garantía \\
%			\midrule
%			Lineal, paso fijo  & $\eta < 1/\lambda_{\max}$
%			& $\rho^k$, $\rho<1$ & Convergencia global \\[3pt]
%			Lineal, paso óptimo & $\eta_k^*$ exacto
%			& $(1-\kappa^{-1})^{2k}$ & Convergencia global \\[3pt]
%			\rowcolor{accentteal!15}
%			$L$-suave (no lineal) & $\eta \leq 1/L$
%			& $O(1/k)$ en $\|\nabla\mathcal{L}\|^2$ & Punto crítico \\[3pt]
%			$L$-suave + convexo  & $\eta \leq 1/L$
%			& $O(1/k)$ en $\mathcal{L}$ & Mínimo global \\
%			\bottomrule
%		\end{tabular}
%	\end{center}
%	
%	\vspace{8pt}
%	\begin{obsbox}{Moraleja para redes neuronales}
%		$\mathcal{L}$ es no convexa en general $\Rightarrow$ solo garantizamos llegar a un
%		\textbf{punto crítico}, no al mínimo global. \\
%		El paso $\eta < 1/L$ es la condición suficiente.
%	\end{obsbox}
%\end{frame}

\begin{frame}{El ciclo completo de entrenamiento}
	\begin{center}
		\begin{tikzpicture}[node distance=7mm and 18mm]
			\node[graphop, fill=accentviolet!15, draw=accentviolet,
			minimum width=38mm] (fwd) {Forward pass};
			\node[graphop, fill=accentteal!15, draw=accentteal,
			minimum width=38mm, below=of fwd] (loss) {Calcular $\mathcal{L}(\theta_k)$};
			\node[graphop, fill=coralbox!15, draw=coralbox,
			minimum width=38mm, below=of loss] (bwd) {Backward pass};
			\node[graphop, fill=accentamber!15, draw=accentamber,
			minimum width=38mm, below=of bwd] (upd) {Actualizar $\theta$};
			
			\draw[thickconn=accentviolet] (fwd) -- node[right,font=\small]{$\hat{f}(x;\theta_k)$} (loss);
			\draw[thickconn=accentteal]   (loss) -- node[right,font=\small]{$\nabla_\theta\mathcal{L}$} (bwd);
			\draw[thickconn=coralbox]     (bwd) -- node[right,font=\small]{$\theta_{k+1}$} (upd);
			
			\draw[thickconn=accentamber, rounded corners=4pt]
			(upd.east) -- ++(8mm,0) -- ++(0,33mm) -- (fwd.east);
			\node[right=12mm of loss, font=\small\color{accentamber}] {repetir};
		\end{tikzpicture}
	\end{center}
	
\end{frame}

\begin{frame}{El ciclo completo de entrenamiento}
	\vspace{2pt}
	\begin{resultbox}{Actualización explícita por capa}
		Para la red $\hat{f} = W_2\,\sigma(W_1 x + \bm{b}_1) + b_2$, un paso de gradiente:
		\[
		W_1 \leftarrow W_1 - \eta\,\frac{\partial\mathcal{L}}{\partial W_1}, \quad
		\bm{b}_1 \leftarrow \bm{b}_1 - \eta\,\frac{\partial\mathcal{L}}{\partial \bm{b}_1}, \quad
		W_2 \leftarrow W_2 - \eta\,\frac{\partial\mathcal{L}}{\partial W_2}, \quad
		b_2 \leftarrow b_2 - \eta\,\frac{\partial\mathcal{L}}{\partial b_2}
		\]
		\vspace{2pt}
		\\
		\textbf{Backprop} calcula los cuatro gradientes en una sola pasada hacia atrás.\\ 	\vspace{2pt} \\
		JAX con \texttt{grad} hace esto automáticamente para cualquier arquitectura.
	\end{resultbox}
	
\end{frame}

% ============================================================
%  SECCIÓN 4: BACKPROPAGATION
% ============================================================
\begin{darkslide}
\begin{frame}[plain]
\begin{tikzpicture}[remember picture, overlay]
  \fill[coralbox,opacity=0.12] (current page.north west) ++(2,-2) circle (3.5cm);
  \fill[accentviolet,opacity=0.10] (current page.south east) ++(-2,2) circle (3cm);
\end{tikzpicture}
\vspace{2cm}
\begin{center}
  {\Large\textcolor{textlight}{\textbf{Parte 4}}}\\[10pt]
  {\large\textcolor{coralbox}{Backpropagation}}\\[6pt]
  {\normalsize\textcolor{textgray}{= autodiff en modo reverso sobre el grafo de la red}}
\end{center}
\end{frame}
\end{darkslide}

\begin{frame}{Regla de la cadena a través de la red}
Recordemos: para la red de 1 neurona,
\[
  z = wx + b, \quad \hat{f} = \sigma(z), \quad \mathcal{L} = (y - \hat{f})^2
\]

\vspace{4pt}
\begin{resultbox}{Grafo computacional}
\begin{center}
\begin{tikzpicture}[node distance=4mm and 16mm]
  \node[graphnode] (x) {$x$};
  \node[graphnode, right=of x] (z) {$z$};
  \node[graphnode, right=of z] (fhat) {$\hat{f}$};
  \node[graphnode, right=of fhat] (L) {$\mathcal{L}$};
  \node[graphnode, above=8mm of z, fill=accent!15, draw=accent] (w) {$w$};
  \node[graphnode, below=8mm of z, fill=accent!15, draw=accent] (b) {$b$};

  \draw[thickconn=accentteal] (x) -- node[above,font=\small]{$\times w$} (z);
  \draw[thickconn=accentteal] (w) -- (z);
  \draw[thickconn=accentteal] (b) -- (z);
  \draw[thickconn=accentteal] (z) -- node[above,font=\small]{$\sigma$} (fhat);
  \draw[thickconn=accentteal] (fhat) -- node[above,font=\small]{$(y-\cdot)^2$} (L);

  \node[below=2mm of x, font=\small\color{textgray}] {dato};
  \node[below=16mm of z, font=\small\color{accent}] {parámetros};
\end{tikzpicture}
\end{center}
\end{resultbox}

\vspace{4pt}
\begin{obsbox}{Regla de la cadena}
\[
  \frac{\partial\mathcal{L}}{\partial w} = \frac{\partial\mathcal{L}}{\partial \hat{f}}\cdot\frac{\partial \hat{f}}{\partial z}\cdot\frac{\partial z}{\partial w}
\]
Calculamos \textbf{de derecha a izquierda} en el grafo — modo reverso.
\end{obsbox}
\end{frame}

\begin{frame}{Backprop: pasada hacia adelante y hacia atrás}
\vspace{-0.5cm}
\begin{columns}[T]
\column{0.48\textwidth}
\begin{resultbox}{Forward pass (valores)}
	\vspace{-0.5cm}
\begin{align*}
  z &= wx + b \\
  \hat{f} &= \sigma(z) \\
  \mathcal{L} &= (y - \hat{f})^2
\end{align*}
Guardamos $z$ y $\hat{f}$ para usarlos en el backward.
\end{resultbox}

\column{0.48\textwidth}
\begin{defbox}{Backward pass (gradientes)}
	\vspace{-0.5cm}
\begin{align*}
  \frac{\partial\mathcal{L}}{\partial \hat{f}} &= -2(y-\hat{f}) \\[4pt]
  \frac{\partial\mathcal{L}}{\partial z} &= \frac{\partial\mathcal{L}}{\partial \hat{f}}\cdot\sigma'(z) \\[4pt]
  \frac{\partial\mathcal{L}}{\partial w} &= \frac{\partial\mathcal{L}}{\partial z}\cdot x \\[4pt]
  \frac{\partial\mathcal{L}}{\partial b} &= \frac{\partial\mathcal{L}}{\partial z}\cdot 1
\end{align*}
\end{defbox}
\end{columns}

\begin{obsbox}{Conexión con la clase anterior}
Esto \textbf{es exactamente} autodiff en modo reverso: una pasada forward, una pasada backward, gradiente completo.
\end{obsbox}
\end{frame}

\begin{frame}{¿Por qué no diferencias finitas?}
\begin{center}
\begin{tabular}{lccc}
\toprule
Método & Exacto & Costo & Escala a $p$ parámetros \\
\midrule
Diferencias finitas & No & $O(p)$ evaluaciones & \textcolor{accent}{$p$ veces más lento} \\
Simbólico & Sí & Explosión combinatoria & \textcolor{accent}{No factible} \\
\rowcolor{accentteal!15}
Autodiff (backprop) & Sí & $O(1)$ pasadas & \textcolor{accentteal}{\textbf{Independiente de $p$}} \\
\bottomrule
\end{tabular}
\end{center}

\vspace{10pt}
\begin{resultbox}{Por qué el gradiente exacto importa}
Una red típica tiene $p \sim 10^6$ parámetros.\\
Con diferencias finitas necesitaríamos $\sim 10^6$ evaluaciones de $\mathcal{L}$ por paso de gradiente.\\
Con backprop: \textbf{2 evaluaciones} (forward + backward).
\end{resultbox}
\end{frame}

% ============================================================
%  EJERCICIO EN GRUPOS
% ============================================================
\begin{darkslide}
\begin{frame}[plain]
\begin{tikzpicture}[remember picture, overlay]
  \fill[coralbox,opacity=0.15] (current page.north east) ++(-1.5,-1.5) circle (3.5cm);
  \fill[accentamber,opacity=0.10] (current page.south west) ++(2,2) circle (3cm);
\end{tikzpicture}
\vspace{2cm}
\begin{center}
  {\Large\textcolor{textlight}{\textbf{Ejercicio}}}\\[10pt]
  {\large\textcolor{coralbox}{Backprop a mano + verificación en JAX}}\\[6pt]
  {\small\textcolor{textgray}{20 minutos}}
\end{center}
\end{frame}
\end{darkslide}

\begin{frame}[fragile]{Ejercicio — Parte A: a mano}
\begin{taskbox}{Red de 2 neuronas, calcular gradientes a mano}
Consideren:
\[
  h_1 = \sigma(w_1 x + b_1),\quad h_2 = \sigma(w_2 x + b_2),\quad \hat{f} = v_1 h_1 + v_2 h_2
\]
\[
  \mathcal{L} = (y - \hat{f})^2, \qquad \sigma(z) = \tanh(z)
\]
Con $x=1$, $y=0.5$, $w_1=0.3$, $b_1=0.1$, $w_2=-0.2$, $b_2=0.4$, $v_1=0.6$, $v_2=0.8$.

\vspace{4pt}
\textbf{(a)} Hagan el forward pass: calculen $h_1$, $h_2$, $\hat{f}$, $\mathcal{L}$.\\[3pt]
\textbf{(b)} Hagan el backward pass: calculen $\partial\mathcal{L}/\partial v_1$, $\partial\mathcal{L}/\partial v_2$, $\partial\mathcal{L}/\partial w_1$, $\partial\mathcal{L}/\partial b_1$.\\[3pt]
%\textbf{(c)} Escriban la regla de la cadena que usaron explícitamente.
\end{taskbox}
\end{frame}

\begin{frame}[fragile]{Ejercicio — Parte B: verificación en JAX}

\begin{lstlisting}[language=Python]
import jax.numpy as jnp
from jax import grad, jit
import jax

def red(params, x, y):
    w1, b1, w2, b2, v1, v2 = params
    h1 = jnp.tanh(w1 * x + b1)
    h2 = jnp.tanh(w2 * x + b2)
    f_hat = v1 * h1 + v2 * h2
    return (y - f_hat) ** 2

params0 = jnp.array([0.3, 0.1, -0.2, 0.4, 0.6, 0.8])
x, y = 1.0, 0.5

# Verificar: coincide con el calculo a mano?
gradientes = grad(red)(params0, x, y)
print(gradientes)

# Un paso de descenso de gradiente con eta = 0.1
eta = 0.1
params1 = params0 - eta * gradientes
print("Loss antes:", red(params0, x, y))
print("Loss despues:", red(params1, x, y))
\end{lstlisting}

\vspace{2pt}
\textcolor{accentteal}{\small \texttt{grad(red)} calcula $\nabla_{\theta}\mathcal{L}$ usando backprop exactamente como lo hicieron a mano.}
\end{frame}

\begin{frame}[fragile]{Ejercicio — Parte C: entrenamiento}

\begin{lstlisting}[language=Python]
import matplotlib.pyplot as plt

# Funcion objetivo: f(x) = sin(2*pi*x)
key = jax.random.PRNGKey(42)
xs = jnp.linspace(0, 1, 50)
ys = jnp.sin(2 * jnp.pi * xs)

def loss_total(params):
    return jnp.mean(jax.vmap(lambda x, y: red(params, x, y))(xs, ys))

grad_loss = jit(grad(loss_total))

# Descenso de gradiente
params = params0
eta = 0.05
historico = []
for _ in range(500):
    params = params - eta * grad_loss(params)
    historico.append(float(loss_total(params)))  # guardar

plt.semilogy(historico)
plt.xlabel("iteracion"); plt.ylabel("L(theta)"); plt.title("Convergencia")
plt.show()
\end{lstlisting}

\vspace{2pt}
\textcolor{accentamber}{\small ¿La red converge? ¿Qué pasa si aumentan $m$ (más neuronas)?}
\end{frame}


% ============================================================
%  SLIDE FINAL
% ============================================================
\begin{darkslide}
\begin{frame}[plain]
\begin{tikzpicture}[remember picture, overlay]
  \fill[accent,opacity=0.15] (current page.north east) ++(-1,-1) circle (3.5cm);
  \fill[accentteal,opacity=0.10] (current page.south west) ++(2,2) circle (2.8cm);
  \fill[accentviolet,opacity=0.08] (current page.north west) ++(1,-2) circle (2cm);
\end{tikzpicture}
\vspace{1.8cm}
\begin{center}
  {\large\textcolor{textgray}{Recapitulando}}\\[14pt]
  {\large\textcolor{textlight}{\textbf{Red neuronal}}}\\
  {\small\textcolor{textgray}{$\hat{f}(x;\theta)$ con bases aprendibles — generaliza todo lo anterior}}\\[10pt]
  {\large\textcolor{textlight}{\textbf{Loss function}}}\\
  {\small\textcolor{textgray}{$\mathcal{L}(\theta)$ mide el error — es la integral que minimizamos}}\\[10pt]
  {\large\textcolor{textlight}{\textbf{Backpropagation}}}\\
  {\small\textcolor{textgray}{autodiff modo reverso sobre el grafo de la red — exacto, $O(1)$ pasadas}}\\[20pt]
  \rule{4cm}{0.4pt}\\[6pt]
  {\small\textcolor{textgray}{\texttt{mbastidaso@unal.edu.co}}}
\end{center}
\end{frame}
\end{darkslide}

\end{document}
